\section{Reading priorities and conclusions to keep provisional}
For the research question, read SuperMimic to understand the source-model
opportunity, PWM to understand models as optimization interfaces, DMO to
separate forward error from backward error, and WM-LOCO to understand the closest
terrain-aware humanoid setting. Then use ContactNets to reason about structure
and $\nabla$Mimic to reason about gradient handling. RWM-U, PETS, MBPO, Dreamer,
and SHAC supply the complementary uncertainty, rollout, representation, and
value-learning principles.

The strongest provisional hypothesis is that a motion-trained model supplies a
useful initialization, but the benefit depends on explicit geometry information
and on preserving useful sensitivities near newly encountered contacts. That
hypothesis has not been demonstrated by the proposed experiments in this review.

The evidence does not justify an a priori interaction count, a guarantee that
structured models require less data, or a universal claim that exact or learned
gradients are better. It does justify experiments that isolate those mechanisms.
The paper-specific formulations and the experimental design should be used as a
map for testing, not as a substitute for new measurements.

\appendix
\section{Symbol crosswalk to the original papers}\label{app:crosswalk}
The mathematical roles stay fixed in this review even when a paper reuses a
letter for a different network. This table is a navigation aid, not a claim that
the complete state representations are equivalent.
\begin{longtable}{@{}L{0.17\linewidth}L{0.37\linewidth}L{0.38\linewidth}@{}}
\toprule Paper & Original role / notation & This review\\\midrule\endhead
PETS & Physical $s$; model $f_\theta$ & Physical $x$; transition $\widehat P_\phi$\\
MBPO & Model $p_\theta$; rollout length $k$ & $\widehat P_\phi$; $h_{\rm roll}$\\
Dreamer & Representation $p_\theta$ & Posterior $\mathsf q_\phi$\\
Dreamer & Transition $q_\theta$ & Prior $\mathsf p_\phi$\\
Dreamer & Action $q_\varphi$; value $v_\psi$ & Actor $\pi_\theta$; value $V_\psi$\\
ContactNets & Geometry parameters $\theta$; impulses $\lambda$ & Geometry/model $\phi$; impulses $j$\\
SHAC & Value $V_\varphi$; target $V_{\varphi'}$ & $V_\psi$; $V_{\bar\psi}$\\
SHAC & Short horizon $h$, full task horizon $H$ & Optimization window $H$; task horizon stated separately\\
PWM & Task embedding $e$ & Task embedding $w$, distinct from terrain encoding $e_t$\\
DMO & State $s$; model $\hat f_\varphi$ & State $x$; model $\Fhat$\\
RWM-U & Uncertainty penalty coefficient $\lambda$ & $\beta_U$, distinct from lambda returns\\
SuperMimic & Model weights $\omega$; survival $c_t$ & Model $\phi$; cumulative survival $C_t$\\
SuperMimic & Reward $r_{t+1}$ & Transition-associated reward $r_t$\\
$\nabla$Mimic & Action cotangent $g_i$; target $y_i$ & $b_i$; physical setpoint target $u_i^{\rm target}$\\
$\nabla$Mimic & Critic $V_\varphi$ & Critic $V_\psi$\\
WM-LOCO & Memory $h_t$; stochastic state $z_t$ & Memory $m_t$; stochastic state $\ell_t$\\
WM-LOCO & Velocity command $c_t$ & Part of task-goal input $g_t$\\
WM-LOCO & Observation feature $e_t$ & $f_t^o$, distinct from environment encoding $e_t$\\
WM-LOCO & Predictive feature projection $g_\psi$ & $\mathcal P_\theta$, assigned to the policy role\\
\bottomrule
\end{longtable}
In particular, $\phi$ never denotes the actor merely because an original paper
uses that letter. A latent state $z$ also never silently becomes the physical
robot state $x$. When a single objective is applied to either, its placeholder
$y$ is explicitly defined.

\section{Primary-source audit and version sensitivity}\label{app:sources}
The accompanying \code{source_manifest.json} records the exact downloaded PDF
URLs, file sizes, and SHA-256 hashes. Source PDFs and copied paper figures are
not redistributed. All figures in this review are synthetic explanations or
original diagrams.

\begin{longtable}{@{}L{0.17\linewidth}L{0.27\linewidth}L{0.48\linewidth}@{}}
\toprule Paper & Reviewed version & Main verification location\\\midrule\endhead
PETS & arXiv 1805.12114v2 & Sections 4--5: model distributions, ensemble fitting, particle planning\\
MBPO & arXiv 1906.08253v3 & Sections 4--5: bounds versus practical short-rollout algorithm\\
Dreamer & arXiv 1912.01603v3 & Sections 2--4: latent components, actor/value returns, reconstruction\\
ContactNets & arXiv 2009.11193v2 & Section 4: learned geometry and impulse-consistency losses\\
SHAC & arXiv 2204.07137v1 & Section 3.2 and Algorithm 1: windows, critic, target update\\
PWM & arXiv 2407.02466v3 & Section 3.3: pretraining and first-order policy extraction\\
DMO & CoRL 2025 PMLR PDF & Section 3 and Appendix A.1.3: forward/backward decoupling\\
RWM-U & arXiv 2504.16680v3 & Sections 3--4: recurrent predictions, uncertainty, MOPO-PPO\\
SuperMimic & arXiv 2610.28480v1 & Sections 3.1--3.3: online fitting, context, soft-survival imagination\\
$\nabla$Mimic & arXiv 2610.41654v1 & Sections 4--6 and 8: clipping, targets, derivative semantics, limitations\\
WM-LOCO & arXiv 2609.02542v2 & Section II: RSSM and PPO conditioning; terrain/reward evaluation\\
\bottomrule
\end{longtable}

RWM-U's reviewed 2026 version revises its 2025 preprint. WM-LOCO's second version
also changes the experimental presentation relative to its first version. Claims
must be read from the cited version rather than assembled from incompatible
abstracts or web summaries.

For reproduction, check implementation details not resolved by a compact paper
equation: normalization, unit conversions, terminal bootstrapping, uncertainty
scalarization, gradient clipping placement, and whether a learned feature is
detached during a policy update. A sum of losses does not by itself specify
every module's gradient routing.

\section{Unified algorithm templates}\label{app:algorithms}
These templates are review-level reconstructions. They express the different
information and derivative paths, not line-for-line reproductions of source code.

\subsection{Fully imagined pathwise policy update}
\begin{enumerate}
\item Sample real replay starts and infer any required latent/history state.
\item Hold model weights fixed. For $H$ steps, construct the current task and
environment inputs, evaluate the actor, predict the successor, and compute reward
and continuation. Advance reference and environment queries consistently.
\item Add a value tail; backpropagate into actor parameters through model
operations and state-dependent inputs.
\item Detach the resulting value targets and fit the critic. Fit the model in
its own supervised phase using real transitions.
\end{enumerate}
The model is fixed within the actor update, not necessarily throughout training.
The critic loss being small does not validate the imagined dynamics.

\subsection{Simulator-forward, learned-model-backward update}
\begin{enumerate}
\item Run the current policy in the physical simulator and record actual states.
\item At each visited input, evaluate the learned transition and install a
forward-value correction such as Equation~\eqref{eq:dmo}.
\item Differentiate the simulator-valued return using those learned local
derivatives. Train the critic on simulator-valued rewards and states.
\item Count the simulator trajectory transitions in the interaction budget;
update the model from executed transitions separately.
\end{enumerate}
This removes imagined forward drift, not local derivative error.

\subsection{Predictive-feature PPO update}
\begin{enumerate}
\item Infer predictive recurrent memory from actual observations and past actions.
\item Give a projected memory feature to the actor and collect simulator rollouts.
\item Compute PPO advantages and the clipped policy objective from those rollouts.
\item Train predictive heads and latent dynamics on sequences with their model
loss. Verify cross-module gradient routing in the chosen implementation.
\end{enumerate}
The presence of a predictive latent simulator does not make this a pathwise
policy-gradient method through imagined transitions.

\subsection{Planning or synthetic-data control}
PETS-style control evaluates candidate action sequences under model uncertainty,
executes one action, and replans. MBPO-style learning creates short synthetic
transition branches and uses them in a separate actor--critic learner. Both use
forward predictions without requiring a policy gradient through a chain of
transition Jacobians.

\section{A worksheet for reading future papers}
For a new paper, fill in the following before comparing its score with ours:
\begin{enumerate}
\item What information is supplied to the model, actor, and critic separately?
Does the geometry exist in those inputs?
\item Is the predicted state physical, latent, or an observation history? Is it
sufficient for the task, or only a useful approximation?
\item What advances a rollout: physical simulation, a transition mean, samples,
or an observation-conditioned posterior?
\item What exactly is differentiated to update the actor? Where are gradients
stopped, clipped, replaced, or approximated?
\item What supervises each predictive head? Are analytic rewards distinguished
from learned rewards?
\item Which parameters are reused, frozen, or adapted across environments?
What target interaction is actually counted?
\item Does the evaluation test complete tasks and held-out contact geometry,
or only training instances and short prefixes?
\item Can architecture size, reward shaping, privileged information, or a
pretrained policy explain the apparent model benefit?
\end{enumerate}
This worksheet preserves the common mathematical comparison even when future
papers use new terminology for familiar components.
